\chapter{Явный унитарный представитель обменного класса KO7}

Абстрактный символ $\xi_{15}$ теперь проверяется независимо в унитарной
модели Бурсема--Лоринга. В этой модели $KO_7=KO_{-1}$ представлен
унитариями $V$, удовлетворяющими $V^\tau=V$; см. определение 6.10,
теорему 7.1 и таблицу 3 в \path{arXiv:1504.03284v3}.

Пусть $q_0\in M_{105}(\mathbb C)$ --- проектор ранга пятнадцать. После
комплексификации обменной вещественной алгебры получаются две ветви, и
положим
\begin{equation}
 V_{15}(z)=\left(
 zq_0+1-q_0,
 z^{-1}\overline q_0+1-\overline q_0
 \right).
 \label{eq:v5-ko7-explicit-unitary}
\end{equation}
Обе компоненты унитарны. Линейная антимультипликативная Real-инволюция
имеет вид
\begin{equation}
 \tau(f_+,f_-)(z)=
 \bigl(f_-(\bar z)^T,f_+(\bar z)^T\bigr).
\end{equation}
Так как $q_0^T=\overline q_0$, непосредственно получаем
\begin{equation}
 V_{15}^{\tau}=V_{15}.
\end{equation}
Следовательно, \eqref{eq:v5-ko7-explicit-unitary} действительно задаёт
класс $KO_7$, а не только пару комплексных унитариев.

Комплексные winding-классы двух компонент равны
\begin{equation}
 ([V_+],[V_-])=(15,-15).
\end{equation}
Компрессия проектором Харди даёт
\begin{equation}
 \operatorname{ind}T_{V_+}=-15,
 \qquad
 \operatorname{ind}T_{V_-}=+15.
\end{equation}
Это точно совпадает с $c_6(15)=(-15,+15)$.

\begin{remark}[Отброшенное удвоение]
Формула $pu+(1-p)u^*$ выглядит естественно, но представляет
$2[p]-[1]$. На ранге пятнадцать она дала бы абсолютный класс тридцать.
Она не используется. Разнесение $z$ и $z^{-1}$ по двум Real-обменным
ветвям сохраняет класс пятнадцать без удвоения.
\end{remark}

\begin{theorem}[Явный представитель]
Унитарий $V_{15}$ представляет класс $\xi_{15}\in KO_7$, а его
вещественная Toeplitz boundary равна классу пятнадцать в $KO_6$ с
точностью до общей ориентации. Нормированный вес равен $15/105=1/7$.
\end{theorem}

Теперь замкнуты группа, степень, Real-симметрия, явный унитарий,
неограниченный extension-оператор и boundary-образ. Не получены
пространственное действие, конечная энергия, масса и масштаб.

\begin{protocol}
\begin{enumerate}
 \item размер ветви: \textbf{105};
 \item ранг $q_0$: \textbf{15};
 \item условие $V^\tau=V$: \textbf{выполнено};
 \item complexification: \textbf{$(15,-15)$};
 \item Toeplitz boundary: \textbf{$(-15,+15)$};
 \item наивное удвоение до 30: \textbf{исключено};
 \item вещественный класс: \textbf{15 с точностью до ориентации};
 \item вес: \textbf{$1/7$};
 \item явный KO7-представитель: \textbf{построен};
 \item физическое действие: \textbf{не получено};
 \item следующий гейт: \textbf{вариационный разрыв родительского действия}.
\end{enumerate}
\end{protocol}

Машинный сертификат сохранён в
\path{s2t_v5_real_toeplitz_ko7_unitary_representative_gate_results.json}.